%% Sample Exam 2 -- Medical Image Analysis (ETH Zurich, 2026)
%%
%% Written in the style of Example_Exam_with_Answers.pdf.
%%
%% Compile the version WITH solutions:      pdflatex "\def\withanswers{}%% Sample Exam 2 -- Medical Image Analysis (ETH Zurich, 2026)
%%
%% Written in the style of Example_Exam_with_Answers.pdf.
%%
%% Compile the version WITH solutions:      pdflatex "\def\withanswers{}%% Sample Exam 2 -- Medical Image Analysis (ETH Zurich, 2026)
%%
%% Written in the style of Example_Exam_with_Answers.pdf.
%%
%% Compile the version WITH solutions:      pdflatex "\def\withanswers{}%% Sample Exam 2 -- Medical Image Analysis (ETH Zurich, 2026)
%%
%% Written in the style of Example_Exam_with_Answers.pdf.
%%
%% Compile the version WITH solutions:      pdflatex "\def\withanswers{}\input{Sample_Exam_2.tex}"
%% Compile the blank version (no answers):  pdflatex Sample_Exam_2.tex
%%
\ifx\withanswers\undefined
  \documentclass[11pt]{exam}
\else
  \documentclass[11pt,answers]{exam}
\fi

\usepackage[margin=2.4cm]{geometry}
\usepackage{amsmath,amssymb}
\usepackage{array}
\usepackage[T1]{fontenc}

% ---------------------------------------------------------------- page style
\pagestyle{foot}
\firstpageheader{}{}{}
\firstpagefooter{}{}{}
\runningheader{}{}{}
\runningfooter{}{}{Page \thepage}

% -------------------------------------------------------------- multiplechoice
% Reproduce the "(A) (B) (C) ..." labelling of the original exam.
\renewcommand{\choicelabel}{$\bigcirc$~(\Alph{choice})}

% ---------------------------------------------------------- true/false macros
% \TFT -> the statement is TRUE, \TFF -> the statement is FALSE.
% In the answer version the correct word is set in capitals and bold.
\newcommand{\TFT}{\ifprintanswers\textbf{TRUE}\hspace{1.1em}False\else True\hspace{1.1em}False\fi}
\newcommand{\TFF}{True\hspace{1.1em}\ifprintanswers\textbf{FALSE}\else False\fi}

% ------------------------------------------------------------- writing space
% Reserves room to write in the blank version; collapses in the answer version.
\newcommand{\answerspace}[1]{\ifprintanswers\else\vspace{#1}\fi}

\begin{document}

% ------------------------------------------------------------------ title page
\begin{center}
  {\LARGE\bfseries Sample Exam}\\[2mm]
  {\Large Medical Image Analysis}\\[6mm]
\end{center}

\vspace{4mm}
\noindent Student name:\hrulefill

\vspace{6mm}
\noindent Immatriculation number:\hrulefill

\vspace{10mm}

\begin{center}
  {\large\bfseries Exam Questions}
\end{center}

\noindent\textbf{Note:} All the questions have equal weight.

\vspace{4mm}

\begin{questions}

% =============================================================== Question 1 ===
\question
In 2D acquisitions, 2D image slices are acquired in one pre-defined direction one after the
other, and the final 3D volume is created by stacking the 2D slices. For instance, in axial
acquisitions 2D slices are acquired in the transverse plane, in sagittal acquisitions in the
sagittal plane, and in coronal acquisitions in the coronal plane. The voxel sizes of 2D
acquisitions are given based on the acquisition direction, as in-plane and through-plane voxel
sizes, where the first is represented with two values (pixel size in the imaging plane) and the
latter corresponds to a single value (length of the voxel in the third dimension). Recall that
the field-of-view is the real-world size of the entire image, i.e.\ image size $\times$ voxel
size. Based on this please compute the following.

\begin{parts}

  \part A \emph{sagittal} acquisition with voxel size in-plane $0.5 \times 0.5$~mm$^2$ with 384
  pixels in both directions and voxel size through-plane $4.0$~mm with 60 pixels is acquired.
  What is the field-of-view in the head-to-toe, left-to-right and anterior-to-posterior
  directions, respectively?

  \begin{solution}
    Answer: $192 \times 240 \times 192$~mm$^3$.

    The sagittal plane contains the head-to-toe and the anterior-to-posterior directions, so
    both in-plane directions give $384 \times 0.5 = 192$~mm. The through-plane direction of a
    sagittal acquisition is left-to-right: $60 \times 4.0 = 240$~mm.
  \end{solution}
\answerspace{2.2cm}

  \part One year later the same patient is scanned again, this time with an \emph{axial}
  acquisition with voxel size in-plane $0.8 \times 0.8$~mm$^2$ with a field-of-view of 192~mm in
  both directions, and voxel size through-plane of $2.0$~mm with a field-of-view of 240~mm.
  What are the number of voxels in the head-to-toe, left-to-right and anterior-to-posterior
  directions, respectively?

  \begin{solution}
    Answer: $120 \times 240 \times 240$~voxels.

    The through-plane direction of an axial acquisition is head-to-toe: $240 / 2.0 = 120$. Both
    in-plane directions give $192 / 0.8 = 240$.
  \end{solution}
\answerspace{2.2cm}

  \part You want to compare a volumetric measurement of the same anatomical structure in the two
  images. Following the guideline given in the lecture, you resample the second image (b) onto
  the voxel grid of the first image (a). What are the resize factors along the head-to-toe,
  left-to-right and anterior-to-posterior directions, and what is the resulting image size?
  Why do you match the voxel size and not the image size?

  \begin{solution}
    Answer: resize factors $4.0 \times 0.2 \times 1.6$, resulting image size
    $480 \times 48 \times 384$~voxels.

    The resize factor is the ratio of the old to the new voxel size, per direction:
    head-to-toe $2.0/0.5 = 4.0$, left-to-right $0.8/4.0 = 0.2$, anterior-to-posterior
    $0.8/0.5 = 1.6$. Applying these to (b): $120 \cdot 4.0 = 480$, $240 \cdot 0.2 = 48$,
    $240 \cdot 1.6 = 384$. The field-of-view is unchanged ($240 \times 192 \times 192$~mm$^3$),
    as it must be.

    Measurements taken in an image correspond to real-world quantities only through the voxel
    size, so the voxel size is what has to be matched. Matching the image sizes instead would
    silently change the effective voxel size (and here also the aspect ratio), and the two
    measurements would no longer be comparable.
  \end{solution}
\answerspace{4.5cm}

\end{parts}

% =============================================================== Question 2 ===
\question
Many noise suppression methods can be written as a weighted average
$\hat I(x) = \sum_{y} w(x,y)\, J(y)$ of the noisy image $J$, with $w(x,y) \geq 0$ and
$\sum_y w(x,y) = 1$. Which of the two weight definitions below is the one used by the
\emph{non-local means} filter? $N(x)$ denotes the image patch centred at $x$, $Z(x)$ is a
normalisation constant and $h$ the degree of filtering. Please circle the correct choice.

\vspace{2mm}
\begin{center}
  $\displaystyle w(x,y) = \frac{1}{Z(x)} \exp\left(-\frac{\|J(x) - J(y)\|_2^2}{h^2}\right)$
  \hspace{1.6cm}
  $\displaystyle w(x,y) = \frac{1}{Z(x)} \exp\left(-\frac{\|J(N(x)) - J(N(y))\|_2^2}{h^2}\right)$
\end{center}
\vspace{2mm}

\begin{solution}
  Answer: Right one.

  Non-local means compares the whole geometric configuration of a \emph{patch}, not a single
  grey value. The left expression compares single (noisy) intensities and is the
  Yaroslavsky / bilateral type of neighbourhood filter, which is much less robust: two pixels
  can share an intensity by chance while sitting in completely different structures.
\end{solution}
\answerspace{1.2cm}

% =============================================================== Question 3 ===
\question
Deformable intensity-based registration is usually written in one of the two equivalent forms
below. The first is derived from a probabilistic model, the second is energy-based;
$T_\theta$ is the transformation with parameters $\theta$, $I$ the target and $J$ the moving
image.

\vspace{2mm}
\begin{center}
  $\displaystyle \max_\theta\ \log p(I \mid T_\theta \circ J) + \log p(T_\theta)$
  \hspace{1.8cm}
  $\displaystyle \min_\theta\ d(I, T_\theta \circ J) + \lambda\, r(T_\theta^{-1})$
\end{center}
\vspace{2mm}

\begin{parts}
  \part Please indicate which term corresponds to which term in the two formulations.

  \begin{solution}
    Answer: $\log p(I \mid T_\theta \circ J) \Longleftrightarrow d(I, T_\theta \circ J)$ and
    $\log p(T_\theta) \Longleftrightarrow \lambda\, r(T_\theta^{-1})$.
  \end{solution}
\answerspace{2.2cm}

  \part Deformable registration is an ill-posed problem. Which of the two terms addresses this,
  and how?

  \begin{solution}
    Answer: the regularisation term $\lambda\, r(T_\theta^{-1})$, equivalently the prior
    $p(T_\theta)$.

    Many very different transformations explain the intensities about equally well, so the data
    term alone does not determine a unique solution. The regulariser expresses a preference among
    them --- e.g.\ $\ell_2$ regularisation of $\nabla u$ prefers smoothly varying displacement
    fields, $\ell_1$ prefers piecewise constant ones, and elastic / hyperelastic energies prefer
    small or locally rigid deformations.
  \end{solution}
\answerspace{3.5cm}
\end{parts}

% =============================================================== Question 4 ===
\question
Please indicate whether the statement is true or false by circling the correct choice in each of
the following.

\begin{parts}
  \part \TFF\ \ In the discrete setting, the \emph{forward} transformation is preferred over
  the backward (inverse) transformation for resampling an image, because it guarantees that
  every pixel of the output grid is assigned a value.

  \part \TFT\ \ Nearest-neighbour interpolation is the appropriate choice for transforming a
  label (segmentation) map, because it does not introduce new values.

  \part \TFF\ \ Bilinear interpolation produces a function that is both continuous and
  differentiable.

  \part \TFF\ \ In 3D, Keys' cubic convolution (tricubic) interpolation uses only the 8 corner
  voxels of the cell that contains the query point.
\end{parts}

\begin{solution}
  (a) FALSE --- it is exactly the other way around. Under a forward mapping the transformed
  input pixels do not land on output grid points, two input pixels may land on the same output
  pixel, and some output pixels receive nothing at all (holes). The backward mapping asks, for
  each output pixel, where it comes from in the input image, and interpolates there; every
  output pixel therefore gets a value.

  (b) TRUE --- nearest neighbour only copies existing values, so no labels that do not exist can
  appear. Bilinear or bicubic interpolation would create spurious intermediate label values at
  the boundaries.

  (c) FALSE --- bilinear interpolation yields a piecewise-linear function: continuous, but not
  differentiable at the grid lines. Bicubic interpolation is the one that is both continuous and
  differentiable.

  (d) FALSE --- the 8 corner voxels are what \emph{trilinear} interpolation uses. Keys' kernel
  has support $[-2,2]$, covering 4 samples per dimension, i.e.\ $4\times4 = 16$ pixels in 2D and
  $4\times4\times4 = 64$ voxels in 3D.
\end{solution}

% =============================================================== Question 5 ===
\question
Mixture models and Expectation-Maximization (EM) segmentation are essential building blocks for
many segmentation algorithms. Consider a multi-modal image in which every pixel carries
\emph{two} intensity values (for instance a co-registered T1-weighted and T2-weighted pair). How
many parameters does a Gaussian Mixture Model (GMM) with 3 components have? Count, for each
component, the mean vector, the full covariance matrix and the mixture weight.

\begin{solution}
  Answer: $3 \times (2 + 3 + 1) = 18$.

  Per component: a mean vector in $\mathbb{R}^2$ gives 2 parameters; a full $2\times2$ covariance
  matrix is symmetric and therefore has 3 free parameters ($\sigma_1^2$, $\sigma_2^2$ and one
  covariance); the mixture weight gives 1. That is 6 parameters per component and 18 in total.
\end{solution}
\answerspace{3.0cm}

% =============================================================== Question 6 ===
\question
An Active Shape Model (ASM) is fitted to an image iteratively. Once the model has been
initialised and candidate points have been found by sampling image features along the surface
normals of the model points, each iteration updates the model in two parts. Please briefly
describe what these two parts do, and state what keeps the fitted shape anatomically plausible.
You do not need to write down the equations.

\begin{solution}
  Answer: \emph{Pose estimation} computes a global transformation --- e.g.\ a similarity
  transform, i.e.\ translation, rotation and isotropic scaling --- that best aligns the model
  points with the candidate points found along the normals, for instance by minimising the mean
  Euclidean distance between the two point sets (ICP). \emph{Shape estimation} then removes the
  pose and computes the shape parameters $b$ of the PCA (point distribution) model in closed
  form, which deforms the mean shape along its modes of variation. The two are iterated until
  convergence.

  Plausibility is enforced by the model itself: only the leading principal modes are kept (e.g.\
  as many as needed to explain about 90\% of the variance), and each shape parameter is
  constrained to $|b_i| < 3\sigma_i$, so the fit can only produce shapes that are typical of the
  training population.
\end{solution}
\answerspace{7.5cm}

% =============================================================== Question 7 ===
\question
Thin-plate splines (TPS) and free-form deformations (FFD) with cubic B-splines are both
interpolation-based models for non-linear transformations. Which of the following statements on
their similarities and differences is \textbf{inaccurate}?

\begin{choices}
  \choice Both models place displacement vectors on a set of control points and interpolate in
  between with a smooth basis function, so that the resulting transformation is smooth for any
  value of the parameters.
  \choice The influence of a single control point is global in TPS, whereas for FFD with cubic
  B-splines it is local, because the B-spline basis functions have finite support.
  \choice In FFD, increasing the number of control points, i.e.\ using a finer control grid,
  results in smoother transformations.
  \choice For a given set of $N$ control points in 3D, both models are parameterised by the
  $3N$ coefficients of the displacement vectors at the control points.
\end{choices}

\begin{solution}
  Answer: C.

  It is the other way around: a finer control grid gives the model more degrees of freedom
  locally and therefore produces \emph{less} smooth transformations. This is precisely what makes
  multi-resolution registration possible --- start with few control points for a coarse, very
  smooth alignment and progressively refine the grid.
\end{solution}

% =============================================================== Question 8 ===
\question
Multi-atlas and patch-based methods were the state of the art in segmentation for a long time.
Which of the following statements about them is \textbf{incorrect}?

\begin{choices}
  \choice In multi-atlas segmentation, every atlas is registered to the unseen image, each atlas
  produces its own propagated segmentation, and these are combined in a decision fusion step such
  as majority voting or STAPLE.
  \choice STAPLE estimates the most probable underlying true segmentation and, simultaneously,
  the sensitivity and specificity of each of the input segmentations; it is solved iteratively
  with an Expectation-Maximization algorithm.
  \choice Selecting only the atlases that are similar to the unseen image before fusing performs
  worse than naively fusing all available atlases, because discarding atlases throws away
  information.
  \choice Patch-based label fusion requires only a coarse, rigid or affine registration, because
  correspondence is re-established locally by comparing patches inside a search window.
\end{choices}

\begin{solution}
  Answer: C.

  Atlas selection followed by fusion \emph{outperforms} naive fusion of all atlases: atlases that
  are dissimilar to the target contribute mostly registration errors, and removing them improves
  the fused result. The same reasoning underlies weighted voting, where atlases are weighted by
  how well they match the target rather than being discarded outright.
\end{solution}

% =============================================================== Question 9 ===
\question
Convolutional neural networks and Vision Transformers (ViT) build very different assumptions
into the model. Which of the following statements is \textbf{incorrect}?

\begin{choices}
  \choice In CNNs, locality and the two-dimensional neighbourhood structure of the image act as
  strong inductive biases for vision tasks.
  \choice In a ViT, self-attention is global and only the MLP layers act locally, so that
  splitting the image into patches is essentially the only two-dimensional inductive bias built
  into the model.
  \choice Because a ViT builds in fewer assumptions about images, it typically requires
  \emph{less} training data than a comparable CNN to reach the same accuracy.
  \choice Without positional encodings, a transformer encoder is equivariant to permutations of
  its input tokens, which is why positional information has to be added to the token embeddings.
\end{choices}

\begin{solution}
  Answer: C.

  Fewer built-in assumptions means the geometric properties of images that a CNN gets for free
  have to be learned from the data instead. ViTs therefore generally require \emph{more} training
  data than a comparable CNN, not less.
\end{solution}

% ============================================================== Question 10 ===
\question
You are training a U-Net to segment a small lesion that occupies far less than 1\% of the image.
Which of the following statements is correct?

\begin{choices}
  \choice With pixel-wise cross-entropy, the loss contributed by the lesion is typically much
  smaller than the loss contributed by the rest of the image, so the lesion can be dominated
  during training.
  \choice The S\o{}rensen-Dice coefficient computed on the binary label sets is differentiable
  and can therefore be used directly as a training loss.
  \choice The Dice loss makes the contribution of a class proportional to its size, which is why
  it helps for small structures.
  \choice The skip connections of the U-Net remove any need to reconsider the loss function,
  because they restore the resolution lost in the encoder.
  \choice Multi-class cross-entropy cannot be used when more than two classes are present.
\end{choices}

\begin{solution}
  Answer: A.

  The cross-entropy loss decomposes into a sum over pixels,
  $\mathcal{L} = \mathcal{L}_{\Omega_{small}} + \mathcal{L}_{\overline{\Omega_{small}}}$, and
  since $|\Omega_{small}| \ll |\overline{\Omega_{small}}|$ the first term is normally negligible.
  (B) is wrong because the Dice coefficient as a set measure is \emph{not} differentiable, which
  is why the soft Dice loss is used instead. (C) is wrong in the opposite direction: the soft Dice
  loss gives each class a contribution of comparable magnitude \emph{regardless} of its size,
  which is exactly why it helps here. (D) confuses resolution with class imbalance --- skip
  connections address the former, not the latter.
\end{solution}

% ============================================================== Question 11 ===
\question
Consider the encoder of a segmentation network that takes grey-scale input images of size
$128 \times 128$ with one input channel and has the following architecture:

\begin{itemize}
  \item \textbf{Conv1}: 32 kernels with kernel size $K = 3 \times 3$, stride $S = 1 \times 1$ in
    each direction, and ``same'' padding $P = 1$.
  \item \textbf{Pool1}: Max pooling with $K = 2 \times 2$ and $S = 2$.
  \item \textbf{Conv2}: 64 kernels with $K = 3 \times 3$, $S = 1$ and ``same'' padding $P = 1$.
  \item \textbf{Pool2}: Max pooling with $K = 2 \times 2$ and $S = 2$.
  \item \textbf{Conv3}: 128 kernels with $K = 3 \times 3$, $S = 1$ and ``same'' padding $P = 1$.
  \item \textbf{Pool3}: Max pooling with $K = 2 \times 2$ and $S = 2$.
\end{itemize}

\noindent Circle True or False:

\begin{parts}
  \part \TFT\ \ The global receptive field of a neuron in Pool3 is $22 \times 22$.

  \part \TFF\ \ The global stride of a neuron in Pool3 is $3 \times 3$, one for each pooling
  layer.

  \part \TFT\ \ The total number of convolutional kernel weights (excluding the biases) in
  Conv1, Pool1, Conv2, Pool2, Conv3 and Pool3 is
  $3 \cdot 3 \cdot 1 \cdot 32 + 3 \cdot 3 \cdot 32 \cdot 64 + 3 \cdot 3 \cdot 64 \cdot 128
  = 92\,448$.

  \part \TFF\ \ The output of Pool3 has a spatial size of $16 \times 16$, and each of its
  neurons therefore sees the entire $128 \times 128$ input image.
\end{parts}

\begin{solution}
  (a) TRUE --- track the receptive field $R$ and the jump (global stride) $j$ layer by layer,
  starting from $R = 1$, $j = 1$, using $R \leftarrow R + (K-1)\cdot j$ and $j \leftarrow j \cdot S$:
  Conv1 $\rightarrow R = 3,\ j = 1$; Pool1 $\rightarrow R = 4,\ j = 2$;
  Conv2 $\rightarrow R = 8,\ j = 2$; Pool2 $\rightarrow R = 10,\ j = 4$;
  Conv3 $\rightarrow R = 18,\ j = 4$; Pool3 $\rightarrow R = 22,\ j = 8$.

  (b) FALSE --- strides compound multiplicatively, they do not add: three poolings of stride 2
  give a global stride of $2 \cdot 2 \cdot 2 = 8$, i.e.\ $8 \times 8$.

  (c) TRUE --- $288 + 18\,432 + 73\,728 = 92\,448$. Pooling layers have no learnable parameters
  at all.

  (d) FALSE --- the spatial size is indeed $128 \to 64 \to 32 \to 16$, since ``same'' convolutions
  do not change the size and each pooling halves it. But the output size and the receptive field
  are different things: by (a) a neuron in Pool3 only sees a $22 \times 22$ window of the input,
  far less than the whole image. This is precisely why deeper networks, or skip connections
  combined with more downsampling, are needed to gather global context.
\end{solution}

% ============================================================== Question 12 ===
\question
Which of the following statements on uncertainty estimation, domain shift and learning from
unlabelled examples are correct (multiple answers are possible)?

\begin{choices}
  \choice Standard, deterministic training and inference amount to approximating
  $p(\theta \mid D, M)$ by a Dirac delta at $\theta^*$ and $p(y \mid x, \theta, M)$ by a Dirac
  delta at $f(x;\theta)$, so that the resulting predictive distribution is a point mass carrying
  no uncertainty at all.
  \choice Test-time adaptation still requires labelled samples from the target domain, but fewer
  of them than fine-tuning does.
  \choice Self-training with pseudo-labels works well only if the predictions of the initial
  model are reasonably good; with poor initial estimates it can diverge quickly.
  \choice In the mean-teacher model the teacher is trained by minimising the supervised loss on
  the labelled examples, while the student is the exponential moving average of the teacher.
  \choice Domain generalisation by extreme data augmentation trains only in the source domain,
  but the improvement it brings may be limited to the kinds of variation that the augmentations
  actually cover.
\end{choices}

\begin{solution}
  Answer: (A), (C) and (E) are correct.

  (B) is wrong: test-time adaptation uses \emph{no} labels and no training set at the target
  domain --- it adapts a subset of the parameters at inference time by minimising a label-free
  cost (prediction entropy, an auto-encoder reconstruction loss, or the agreement of the output
  with a source-trained segmentation prior). It is fine-tuning that needs a few labelled target
  samples.

  (D) has the two networks swapped: the student is the one trained on the supervised plus
  consistency loss, and the teacher is the exponential moving average of the student,
  $\phi^i = \alpha \phi^{i-1} + (1-\alpha)\theta^i$.
\end{solution}

\end{questions}

\end{document}
"
%% Compile the blank version (no answers):  pdflatex Sample_Exam_2.tex
%%
\ifx\withanswers\undefined
  \documentclass[11pt]{exam}
\else
  \documentclass[11pt,answers]{exam}
\fi

\usepackage[margin=2.4cm]{geometry}
\usepackage{amsmath,amssymb}
\usepackage{array}
\usepackage[T1]{fontenc}

% ---------------------------------------------------------------- page style
\pagestyle{foot}
\firstpageheader{}{}{}
\firstpagefooter{}{}{}
\runningheader{}{}{}
\runningfooter{}{}{Page \thepage}

% -------------------------------------------------------------- multiplechoice
% Reproduce the "(A) (B) (C) ..." labelling of the original exam.
\renewcommand{\choicelabel}{$\bigcirc$~(\Alph{choice})}

% ---------------------------------------------------------- true/false macros
% \TFT -> the statement is TRUE, \TFF -> the statement is FALSE.
% In the answer version the correct word is set in capitals and bold.
\newcommand{\TFT}{\ifprintanswers\textbf{TRUE}\hspace{1.1em}False\else True\hspace{1.1em}False\fi}
\newcommand{\TFF}{True\hspace{1.1em}\ifprintanswers\textbf{FALSE}\else False\fi}

% ------------------------------------------------------------- writing space
% Reserves room to write in the blank version; collapses in the answer version.
\newcommand{\answerspace}[1]{\ifprintanswers\else\vspace{#1}\fi}

\begin{document}

% ------------------------------------------------------------------ title page
\begin{center}
  {\LARGE\bfseries Sample Exam}\\[2mm]
  {\Large Medical Image Analysis}\\[6mm]
\end{center}

\vspace{4mm}
\noindent Student name:\hrulefill

\vspace{6mm}
\noindent Immatriculation number:\hrulefill

\vspace{10mm}

\begin{center}
  {\large\bfseries Exam Questions}
\end{center}

\noindent\textbf{Note:} All the questions have equal weight.

\vspace{4mm}

\begin{questions}

% =============================================================== Question 1 ===
\question
In 2D acquisitions, 2D image slices are acquired in one pre-defined direction one after the
other, and the final 3D volume is created by stacking the 2D slices. For instance, in axial
acquisitions 2D slices are acquired in the transverse plane, in sagittal acquisitions in the
sagittal plane, and in coronal acquisitions in the coronal plane. The voxel sizes of 2D
acquisitions are given based on the acquisition direction, as in-plane and through-plane voxel
sizes, where the first is represented with two values (pixel size in the imaging plane) and the
latter corresponds to a single value (length of the voxel in the third dimension). Recall that
the field-of-view is the real-world size of the entire image, i.e.\ image size $\times$ voxel
size. Based on this please compute the following.

\begin{parts}

  \part A \emph{sagittal} acquisition with voxel size in-plane $0.5 \times 0.5$~mm$^2$ with 384
  pixels in both directions and voxel size through-plane $4.0$~mm with 60 pixels is acquired.
  What is the field-of-view in the head-to-toe, left-to-right and anterior-to-posterior
  directions, respectively?

  \begin{solution}
    Answer: $192 \times 240 \times 192$~mm$^3$.

    The sagittal plane contains the head-to-toe and the anterior-to-posterior directions, so
    both in-plane directions give $384 \times 0.5 = 192$~mm. The through-plane direction of a
    sagittal acquisition is left-to-right: $60 \times 4.0 = 240$~mm.
  \end{solution}
\answerspace{2.2cm}

  \part One year later the same patient is scanned again, this time with an \emph{axial}
  acquisition with voxel size in-plane $0.8 \times 0.8$~mm$^2$ with a field-of-view of 192~mm in
  both directions, and voxel size through-plane of $2.0$~mm with a field-of-view of 240~mm.
  What are the number of voxels in the head-to-toe, left-to-right and anterior-to-posterior
  directions, respectively?

  \begin{solution}
    Answer: $120 \times 240 \times 240$~voxels.

    The through-plane direction of an axial acquisition is head-to-toe: $240 / 2.0 = 120$. Both
    in-plane directions give $192 / 0.8 = 240$.
  \end{solution}
\answerspace{2.2cm}

  \part You want to compare a volumetric measurement of the same anatomical structure in the two
  images. Following the guideline given in the lecture, you resample the second image (b) onto
  the voxel grid of the first image (a). What are the resize factors along the head-to-toe,
  left-to-right and anterior-to-posterior directions, and what is the resulting image size?
  Why do you match the voxel size and not the image size?

  \begin{solution}
    Answer: resize factors $4.0 \times 0.2 \times 1.6$, resulting image size
    $480 \times 48 \times 384$~voxels.

    The resize factor is the ratio of the old to the new voxel size, per direction:
    head-to-toe $2.0/0.5 = 4.0$, left-to-right $0.8/4.0 = 0.2$, anterior-to-posterior
    $0.8/0.5 = 1.6$. Applying these to (b): $120 \cdot 4.0 = 480$, $240 \cdot 0.2 = 48$,
    $240 \cdot 1.6 = 384$. The field-of-view is unchanged ($240 \times 192 \times 192$~mm$^3$),
    as it must be.

    Measurements taken in an image correspond to real-world quantities only through the voxel
    size, so the voxel size is what has to be matched. Matching the image sizes instead would
    silently change the effective voxel size (and here also the aspect ratio), and the two
    measurements would no longer be comparable.
  \end{solution}
\answerspace{4.5cm}

\end{parts}

% =============================================================== Question 2 ===
\question
Many noise suppression methods can be written as a weighted average
$\hat I(x) = \sum_{y} w(x,y)\, J(y)$ of the noisy image $J$, with $w(x,y) \geq 0$ and
$\sum_y w(x,y) = 1$. Which of the two weight definitions below is the one used by the
\emph{non-local means} filter? $N(x)$ denotes the image patch centred at $x$, $Z(x)$ is a
normalisation constant and $h$ the degree of filtering. Please circle the correct choice.

\vspace{2mm}
\begin{center}
  $\displaystyle w(x,y) = \frac{1}{Z(x)} \exp\left(-\frac{\|J(x) - J(y)\|_2^2}{h^2}\right)$
  \hspace{1.6cm}
  $\displaystyle w(x,y) = \frac{1}{Z(x)} \exp\left(-\frac{\|J(N(x)) - J(N(y))\|_2^2}{h^2}\right)$
\end{center}
\vspace{2mm}

\begin{solution}
  Answer: Right one.

  Non-local means compares the whole geometric configuration of a \emph{patch}, not a single
  grey value. The left expression compares single (noisy) intensities and is the
  Yaroslavsky / bilateral type of neighbourhood filter, which is much less robust: two pixels
  can share an intensity by chance while sitting in completely different structures.
\end{solution}
\answerspace{1.2cm}

% =============================================================== Question 3 ===
\question
Deformable intensity-based registration is usually written in one of the two equivalent forms
below. The first is derived from a probabilistic model, the second is energy-based;
$T_\theta$ is the transformation with parameters $\theta$, $I$ the target and $J$ the moving
image.

\vspace{2mm}
\begin{center}
  $\displaystyle \max_\theta\ \log p(I \mid T_\theta \circ J) + \log p(T_\theta)$
  \hspace{1.8cm}
  $\displaystyle \min_\theta\ d(I, T_\theta \circ J) + \lambda\, r(T_\theta^{-1})$
\end{center}
\vspace{2mm}

\begin{parts}
  \part Please indicate which term corresponds to which term in the two formulations.

  \begin{solution}
    Answer: $\log p(I \mid T_\theta \circ J) \Longleftrightarrow d(I, T_\theta \circ J)$ and
    $\log p(T_\theta) \Longleftrightarrow \lambda\, r(T_\theta^{-1})$.
  \end{solution}
\answerspace{2.2cm}

  \part Deformable registration is an ill-posed problem. Which of the two terms addresses this,
  and how?

  \begin{solution}
    Answer: the regularisation term $\lambda\, r(T_\theta^{-1})$, equivalently the prior
    $p(T_\theta)$.

    Many very different transformations explain the intensities about equally well, so the data
    term alone does not determine a unique solution. The regulariser expresses a preference among
    them --- e.g.\ $\ell_2$ regularisation of $\nabla u$ prefers smoothly varying displacement
    fields, $\ell_1$ prefers piecewise constant ones, and elastic / hyperelastic energies prefer
    small or locally rigid deformations.
  \end{solution}
\answerspace{3.5cm}
\end{parts}

% =============================================================== Question 4 ===
\question
Please indicate whether the statement is true or false by circling the correct choice in each of
the following.

\begin{parts}
  \part \TFF\ \ In the discrete setting, the \emph{forward} transformation is preferred over
  the backward (inverse) transformation for resampling an image, because it guarantees that
  every pixel of the output grid is assigned a value.

  \part \TFT\ \ Nearest-neighbour interpolation is the appropriate choice for transforming a
  label (segmentation) map, because it does not introduce new values.

  \part \TFF\ \ Bilinear interpolation produces a function that is both continuous and
  differentiable.

  \part \TFF\ \ In 3D, Keys' cubic convolution (tricubic) interpolation uses only the 8 corner
  voxels of the cell that contains the query point.
\end{parts}

\begin{solution}
  (a) FALSE --- it is exactly the other way around. Under a forward mapping the transformed
  input pixels do not land on output grid points, two input pixels may land on the same output
  pixel, and some output pixels receive nothing at all (holes). The backward mapping asks, for
  each output pixel, where it comes from in the input image, and interpolates there; every
  output pixel therefore gets a value.

  (b) TRUE --- nearest neighbour only copies existing values, so no labels that do not exist can
  appear. Bilinear or bicubic interpolation would create spurious intermediate label values at
  the boundaries.

  (c) FALSE --- bilinear interpolation yields a piecewise-linear function: continuous, but not
  differentiable at the grid lines. Bicubic interpolation is the one that is both continuous and
  differentiable.

  (d) FALSE --- the 8 corner voxels are what \emph{trilinear} interpolation uses. Keys' kernel
  has support $[-2,2]$, covering 4 samples per dimension, i.e.\ $4\times4 = 16$ pixels in 2D and
  $4\times4\times4 = 64$ voxels in 3D.
\end{solution}

% =============================================================== Question 5 ===
\question
Mixture models and Expectation-Maximization (EM) segmentation are essential building blocks for
many segmentation algorithms. Consider a multi-modal image in which every pixel carries
\emph{two} intensity values (for instance a co-registered T1-weighted and T2-weighted pair). How
many parameters does a Gaussian Mixture Model (GMM) with 3 components have? Count, for each
component, the mean vector, the full covariance matrix and the mixture weight.

\begin{solution}
  Answer: $3 \times (2 + 3 + 1) = 18$.

  Per component: a mean vector in $\mathbb{R}^2$ gives 2 parameters; a full $2\times2$ covariance
  matrix is symmetric and therefore has 3 free parameters ($\sigma_1^2$, $\sigma_2^2$ and one
  covariance); the mixture weight gives 1. That is 6 parameters per component and 18 in total.
\end{solution}
\answerspace{3.0cm}

% =============================================================== Question 6 ===
\question
An Active Shape Model (ASM) is fitted to an image iteratively. Once the model has been
initialised and candidate points have been found by sampling image features along the surface
normals of the model points, each iteration updates the model in two parts. Please briefly
describe what these two parts do, and state what keeps the fitted shape anatomically plausible.
You do not need to write down the equations.

\begin{solution}
  Answer: \emph{Pose estimation} computes a global transformation --- e.g.\ a similarity
  transform, i.e.\ translation, rotation and isotropic scaling --- that best aligns the model
  points with the candidate points found along the normals, for instance by minimising the mean
  Euclidean distance between the two point sets (ICP). \emph{Shape estimation} then removes the
  pose and computes the shape parameters $b$ of the PCA (point distribution) model in closed
  form, which deforms the mean shape along its modes of variation. The two are iterated until
  convergence.

  Plausibility is enforced by the model itself: only the leading principal modes are kept (e.g.\
  as many as needed to explain about 90\% of the variance), and each shape parameter is
  constrained to $|b_i| < 3\sigma_i$, so the fit can only produce shapes that are typical of the
  training population.
\end{solution}
\answerspace{7.5cm}

% =============================================================== Question 7 ===
\question
Thin-plate splines (TPS) and free-form deformations (FFD) with cubic B-splines are both
interpolation-based models for non-linear transformations. Which of the following statements on
their similarities and differences is \textbf{inaccurate}?

\begin{choices}
  \choice Both models place displacement vectors on a set of control points and interpolate in
  between with a smooth basis function, so that the resulting transformation is smooth for any
  value of the parameters.
  \choice The influence of a single control point is global in TPS, whereas for FFD with cubic
  B-splines it is local, because the B-spline basis functions have finite support.
  \choice In FFD, increasing the number of control points, i.e.\ using a finer control grid,
  results in smoother transformations.
  \choice For a given set of $N$ control points in 3D, both models are parameterised by the
  $3N$ coefficients of the displacement vectors at the control points.
\end{choices}

\begin{solution}
  Answer: C.

  It is the other way around: a finer control grid gives the model more degrees of freedom
  locally and therefore produces \emph{less} smooth transformations. This is precisely what makes
  multi-resolution registration possible --- start with few control points for a coarse, very
  smooth alignment and progressively refine the grid.
\end{solution}

% =============================================================== Question 8 ===
\question
Multi-atlas and patch-based methods were the state of the art in segmentation for a long time.
Which of the following statements about them is \textbf{incorrect}?

\begin{choices}
  \choice In multi-atlas segmentation, every atlas is registered to the unseen image, each atlas
  produces its own propagated segmentation, and these are combined in a decision fusion step such
  as majority voting or STAPLE.
  \choice STAPLE estimates the most probable underlying true segmentation and, simultaneously,
  the sensitivity and specificity of each of the input segmentations; it is solved iteratively
  with an Expectation-Maximization algorithm.
  \choice Selecting only the atlases that are similar to the unseen image before fusing performs
  worse than naively fusing all available atlases, because discarding atlases throws away
  information.
  \choice Patch-based label fusion requires only a coarse, rigid or affine registration, because
  correspondence is re-established locally by comparing patches inside a search window.
\end{choices}

\begin{solution}
  Answer: C.

  Atlas selection followed by fusion \emph{outperforms} naive fusion of all atlases: atlases that
  are dissimilar to the target contribute mostly registration errors, and removing them improves
  the fused result. The same reasoning underlies weighted voting, where atlases are weighted by
  how well they match the target rather than being discarded outright.
\end{solution}

% =============================================================== Question 9 ===
\question
Convolutional neural networks and Vision Transformers (ViT) build very different assumptions
into the model. Which of the following statements is \textbf{incorrect}?

\begin{choices}
  \choice In CNNs, locality and the two-dimensional neighbourhood structure of the image act as
  strong inductive biases for vision tasks.
  \choice In a ViT, self-attention is global and only the MLP layers act locally, so that
  splitting the image into patches is essentially the only two-dimensional inductive bias built
  into the model.
  \choice Because a ViT builds in fewer assumptions about images, it typically requires
  \emph{less} training data than a comparable CNN to reach the same accuracy.
  \choice Without positional encodings, a transformer encoder is equivariant to permutations of
  its input tokens, which is why positional information has to be added to the token embeddings.
\end{choices}

\begin{solution}
  Answer: C.

  Fewer built-in assumptions means the geometric properties of images that a CNN gets for free
  have to be learned from the data instead. ViTs therefore generally require \emph{more} training
  data than a comparable CNN, not less.
\end{solution}

% ============================================================== Question 10 ===
\question
You are training a U-Net to segment a small lesion that occupies far less than 1\% of the image.
Which of the following statements is correct?

\begin{choices}
  \choice With pixel-wise cross-entropy, the loss contributed by the lesion is typically much
  smaller than the loss contributed by the rest of the image, so the lesion can be dominated
  during training.
  \choice The S\o{}rensen-Dice coefficient computed on the binary label sets is differentiable
  and can therefore be used directly as a training loss.
  \choice The Dice loss makes the contribution of a class proportional to its size, which is why
  it helps for small structures.
  \choice The skip connections of the U-Net remove any need to reconsider the loss function,
  because they restore the resolution lost in the encoder.
  \choice Multi-class cross-entropy cannot be used when more than two classes are present.
\end{choices}

\begin{solution}
  Answer: A.

  The cross-entropy loss decomposes into a sum over pixels,
  $\mathcal{L} = \mathcal{L}_{\Omega_{small}} + \mathcal{L}_{\overline{\Omega_{small}}}$, and
  since $|\Omega_{small}| \ll |\overline{\Omega_{small}}|$ the first term is normally negligible.
  (B) is wrong because the Dice coefficient as a set measure is \emph{not} differentiable, which
  is why the soft Dice loss is used instead. (C) is wrong in the opposite direction: the soft Dice
  loss gives each class a contribution of comparable magnitude \emph{regardless} of its size,
  which is exactly why it helps here. (D) confuses resolution with class imbalance --- skip
  connections address the former, not the latter.
\end{solution}

% ============================================================== Question 11 ===
\question
Consider the encoder of a segmentation network that takes grey-scale input images of size
$128 \times 128$ with one input channel and has the following architecture:

\begin{itemize}
  \item \textbf{Conv1}: 32 kernels with kernel size $K = 3 \times 3$, stride $S = 1 \times 1$ in
    each direction, and ``same'' padding $P = 1$.
  \item \textbf{Pool1}: Max pooling with $K = 2 \times 2$ and $S = 2$.
  \item \textbf{Conv2}: 64 kernels with $K = 3 \times 3$, $S = 1$ and ``same'' padding $P = 1$.
  \item \textbf{Pool2}: Max pooling with $K = 2 \times 2$ and $S = 2$.
  \item \textbf{Conv3}: 128 kernels with $K = 3 \times 3$, $S = 1$ and ``same'' padding $P = 1$.
  \item \textbf{Pool3}: Max pooling with $K = 2 \times 2$ and $S = 2$.
\end{itemize}

\noindent Circle True or False:

\begin{parts}
  \part \TFT\ \ The global receptive field of a neuron in Pool3 is $22 \times 22$.

  \part \TFF\ \ The global stride of a neuron in Pool3 is $3 \times 3$, one for each pooling
  layer.

  \part \TFT\ \ The total number of convolutional kernel weights (excluding the biases) in
  Conv1, Pool1, Conv2, Pool2, Conv3 and Pool3 is
  $3 \cdot 3 \cdot 1 \cdot 32 + 3 \cdot 3 \cdot 32 \cdot 64 + 3 \cdot 3 \cdot 64 \cdot 128
  = 92\,448$.

  \part \TFF\ \ The output of Pool3 has a spatial size of $16 \times 16$, and each of its
  neurons therefore sees the entire $128 \times 128$ input image.
\end{parts}

\begin{solution}
  (a) TRUE --- track the receptive field $R$ and the jump (global stride) $j$ layer by layer,
  starting from $R = 1$, $j = 1$, using $R \leftarrow R + (K-1)\cdot j$ and $j \leftarrow j \cdot S$:
  Conv1 $\rightarrow R = 3,\ j = 1$; Pool1 $\rightarrow R = 4,\ j = 2$;
  Conv2 $\rightarrow R = 8,\ j = 2$; Pool2 $\rightarrow R = 10,\ j = 4$;
  Conv3 $\rightarrow R = 18,\ j = 4$; Pool3 $\rightarrow R = 22,\ j = 8$.

  (b) FALSE --- strides compound multiplicatively, they do not add: three poolings of stride 2
  give a global stride of $2 \cdot 2 \cdot 2 = 8$, i.e.\ $8 \times 8$.

  (c) TRUE --- $288 + 18\,432 + 73\,728 = 92\,448$. Pooling layers have no learnable parameters
  at all.

  (d) FALSE --- the spatial size is indeed $128 \to 64 \to 32 \to 16$, since ``same'' convolutions
  do not change the size and each pooling halves it. But the output size and the receptive field
  are different things: by (a) a neuron in Pool3 only sees a $22 \times 22$ window of the input,
  far less than the whole image. This is precisely why deeper networks, or skip connections
  combined with more downsampling, are needed to gather global context.
\end{solution}

% ============================================================== Question 12 ===
\question
Which of the following statements on uncertainty estimation, domain shift and learning from
unlabelled examples are correct (multiple answers are possible)?

\begin{choices}
  \choice Standard, deterministic training and inference amount to approximating
  $p(\theta \mid D, M)$ by a Dirac delta at $\theta^*$ and $p(y \mid x, \theta, M)$ by a Dirac
  delta at $f(x;\theta)$, so that the resulting predictive distribution is a point mass carrying
  no uncertainty at all.
  \choice Test-time adaptation still requires labelled samples from the target domain, but fewer
  of them than fine-tuning does.
  \choice Self-training with pseudo-labels works well only if the predictions of the initial
  model are reasonably good; with poor initial estimates it can diverge quickly.
  \choice In the mean-teacher model the teacher is trained by minimising the supervised loss on
  the labelled examples, while the student is the exponential moving average of the teacher.
  \choice Domain generalisation by extreme data augmentation trains only in the source domain,
  but the improvement it brings may be limited to the kinds of variation that the augmentations
  actually cover.
\end{choices}

\begin{solution}
  Answer: (A), (C) and (E) are correct.

  (B) is wrong: test-time adaptation uses \emph{no} labels and no training set at the target
  domain --- it adapts a subset of the parameters at inference time by minimising a label-free
  cost (prediction entropy, an auto-encoder reconstruction loss, or the agreement of the output
  with a source-trained segmentation prior). It is fine-tuning that needs a few labelled target
  samples.

  (D) has the two networks swapped: the student is the one trained on the supervised plus
  consistency loss, and the teacher is the exponential moving average of the student,
  $\phi^i = \alpha \phi^{i-1} + (1-\alpha)\theta^i$.
\end{solution}

\end{questions}

\end{document}
"
%% Compile the blank version (no answers):  pdflatex Sample_Exam_2.tex
%%
\ifx\withanswers\undefined
  \documentclass[11pt]{exam}
\else
  \documentclass[11pt,answers]{exam}
\fi

\usepackage[margin=2.4cm]{geometry}
\usepackage{amsmath,amssymb}
\usepackage{array}
\usepackage[T1]{fontenc}

% ---------------------------------------------------------------- page style
\pagestyle{foot}
\firstpageheader{}{}{}
\firstpagefooter{}{}{}
\runningheader{}{}{}
\runningfooter{}{}{Page \thepage}

% -------------------------------------------------------------- multiplechoice
% Reproduce the "(A) (B) (C) ..." labelling of the original exam.
\renewcommand{\choicelabel}{$\bigcirc$~(\Alph{choice})}

% ---------------------------------------------------------- true/false macros
% \TFT -> the statement is TRUE, \TFF -> the statement is FALSE.
% In the answer version the correct word is set in capitals and bold.
\newcommand{\TFT}{\ifprintanswers\textbf{TRUE}\hspace{1.1em}False\else True\hspace{1.1em}False\fi}
\newcommand{\TFF}{True\hspace{1.1em}\ifprintanswers\textbf{FALSE}\else False\fi}

% ------------------------------------------------------------- writing space
% Reserves room to write in the blank version; collapses in the answer version.
\newcommand{\answerspace}[1]{\ifprintanswers\else\vspace{#1}\fi}

\begin{document}

% ------------------------------------------------------------------ title page
\begin{center}
  {\LARGE\bfseries Sample Exam}\\[2mm]
  {\Large Medical Image Analysis}\\[6mm]
\end{center}

\vspace{4mm}
\noindent Student name:\hrulefill

\vspace{6mm}
\noindent Immatriculation number:\hrulefill

\vspace{10mm}

\begin{center}
  {\large\bfseries Exam Questions}
\end{center}

\noindent\textbf{Note:} All the questions have equal weight.

\vspace{4mm}

\begin{questions}

% =============================================================== Question 1 ===
\question
In 2D acquisitions, 2D image slices are acquired in one pre-defined direction one after the
other, and the final 3D volume is created by stacking the 2D slices. For instance, in axial
acquisitions 2D slices are acquired in the transverse plane, in sagittal acquisitions in the
sagittal plane, and in coronal acquisitions in the coronal plane. The voxel sizes of 2D
acquisitions are given based on the acquisition direction, as in-plane and through-plane voxel
sizes, where the first is represented with two values (pixel size in the imaging plane) and the
latter corresponds to a single value (length of the voxel in the third dimension). Recall that
the field-of-view is the real-world size of the entire image, i.e.\ image size $\times$ voxel
size. Based on this please compute the following.

\begin{parts}

  \part A \emph{sagittal} acquisition with voxel size in-plane $0.5 \times 0.5$~mm$^2$ with 384
  pixels in both directions and voxel size through-plane $4.0$~mm with 60 pixels is acquired.
  What is the field-of-view in the head-to-toe, left-to-right and anterior-to-posterior
  directions, respectively?

  \begin{solution}
    Answer: $192 \times 240 \times 192$~mm$^3$.

    The sagittal plane contains the head-to-toe and the anterior-to-posterior directions, so
    both in-plane directions give $384 \times 0.5 = 192$~mm. The through-plane direction of a
    sagittal acquisition is left-to-right: $60 \times 4.0 = 240$~mm.
  \end{solution}
\answerspace{2.2cm}

  \part One year later the same patient is scanned again, this time with an \emph{axial}
  acquisition with voxel size in-plane $0.8 \times 0.8$~mm$^2$ with a field-of-view of 192~mm in
  both directions, and voxel size through-plane of $2.0$~mm with a field-of-view of 240~mm.
  What are the number of voxels in the head-to-toe, left-to-right and anterior-to-posterior
  directions, respectively?

  \begin{solution}
    Answer: $120 \times 240 \times 240$~voxels.

    The through-plane direction of an axial acquisition is head-to-toe: $240 / 2.0 = 120$. Both
    in-plane directions give $192 / 0.8 = 240$.
  \end{solution}
\answerspace{2.2cm}

  \part You want to compare a volumetric measurement of the same anatomical structure in the two
  images. Following the guideline given in the lecture, you resample the second image (b) onto
  the voxel grid of the first image (a). What are the resize factors along the head-to-toe,
  left-to-right and anterior-to-posterior directions, and what is the resulting image size?
  Why do you match the voxel size and not the image size?

  \begin{solution}
    Answer: resize factors $4.0 \times 0.2 \times 1.6$, resulting image size
    $480 \times 48 \times 384$~voxels.

    The resize factor is the ratio of the old to the new voxel size, per direction:
    head-to-toe $2.0/0.5 = 4.0$, left-to-right $0.8/4.0 = 0.2$, anterior-to-posterior
    $0.8/0.5 = 1.6$. Applying these to (b): $120 \cdot 4.0 = 480$, $240 \cdot 0.2 = 48$,
    $240 \cdot 1.6 = 384$. The field-of-view is unchanged ($240 \times 192 \times 192$~mm$^3$),
    as it must be.

    Measurements taken in an image correspond to real-world quantities only through the voxel
    size, so the voxel size is what has to be matched. Matching the image sizes instead would
    silently change the effective voxel size (and here also the aspect ratio), and the two
    measurements would no longer be comparable.
  \end{solution}
\answerspace{4.5cm}

\end{parts}

% =============================================================== Question 2 ===
\question
Many noise suppression methods can be written as a weighted average
$\hat I(x) = \sum_{y} w(x,y)\, J(y)$ of the noisy image $J$, with $w(x,y) \geq 0$ and
$\sum_y w(x,y) = 1$. Which of the two weight definitions below is the one used by the
\emph{non-local means} filter? $N(x)$ denotes the image patch centred at $x$, $Z(x)$ is a
normalisation constant and $h$ the degree of filtering. Please circle the correct choice.

\vspace{2mm}
\begin{center}
  $\displaystyle w(x,y) = \frac{1}{Z(x)} \exp\left(-\frac{\|J(x) - J(y)\|_2^2}{h^2}\right)$
  \hspace{1.6cm}
  $\displaystyle w(x,y) = \frac{1}{Z(x)} \exp\left(-\frac{\|J(N(x)) - J(N(y))\|_2^2}{h^2}\right)$
\end{center}
\vspace{2mm}

\begin{solution}
  Answer: Right one.

  Non-local means compares the whole geometric configuration of a \emph{patch}, not a single
  grey value. The left expression compares single (noisy) intensities and is the
  Yaroslavsky / bilateral type of neighbourhood filter, which is much less robust: two pixels
  can share an intensity by chance while sitting in completely different structures.
\end{solution}
\answerspace{1.2cm}

% =============================================================== Question 3 ===
\question
Deformable intensity-based registration is usually written in one of the two equivalent forms
below. The first is derived from a probabilistic model, the second is energy-based;
$T_\theta$ is the transformation with parameters $\theta$, $I$ the target and $J$ the moving
image.

\vspace{2mm}
\begin{center}
  $\displaystyle \max_\theta\ \log p(I \mid T_\theta \circ J) + \log p(T_\theta)$
  \hspace{1.8cm}
  $\displaystyle \min_\theta\ d(I, T_\theta \circ J) + \lambda\, r(T_\theta^{-1})$
\end{center}
\vspace{2mm}

\begin{parts}
  \part Please indicate which term corresponds to which term in the two formulations.

  \begin{solution}
    Answer: $\log p(I \mid T_\theta \circ J) \Longleftrightarrow d(I, T_\theta \circ J)$ and
    $\log p(T_\theta) \Longleftrightarrow \lambda\, r(T_\theta^{-1})$.
  \end{solution}
\answerspace{2.2cm}

  \part Deformable registration is an ill-posed problem. Which of the two terms addresses this,
  and how?

  \begin{solution}
    Answer: the regularisation term $\lambda\, r(T_\theta^{-1})$, equivalently the prior
    $p(T_\theta)$.

    Many very different transformations explain the intensities about equally well, so the data
    term alone does not determine a unique solution. The regulariser expresses a preference among
    them --- e.g.\ $\ell_2$ regularisation of $\nabla u$ prefers smoothly varying displacement
    fields, $\ell_1$ prefers piecewise constant ones, and elastic / hyperelastic energies prefer
    small or locally rigid deformations.
  \end{solution}
\answerspace{3.5cm}
\end{parts}

% =============================================================== Question 4 ===
\question
Please indicate whether the statement is true or false by circling the correct choice in each of
the following.

\begin{parts}
  \part \TFF\ \ In the discrete setting, the \emph{forward} transformation is preferred over
  the backward (inverse) transformation for resampling an image, because it guarantees that
  every pixel of the output grid is assigned a value.

  \part \TFT\ \ Nearest-neighbour interpolation is the appropriate choice for transforming a
  label (segmentation) map, because it does not introduce new values.

  \part \TFF\ \ Bilinear interpolation produces a function that is both continuous and
  differentiable.

  \part \TFF\ \ In 3D, Keys' cubic convolution (tricubic) interpolation uses only the 8 corner
  voxels of the cell that contains the query point.
\end{parts}

\begin{solution}
  (a) FALSE --- it is exactly the other way around. Under a forward mapping the transformed
  input pixels do not land on output grid points, two input pixels may land on the same output
  pixel, and some output pixels receive nothing at all (holes). The backward mapping asks, for
  each output pixel, where it comes from in the input image, and interpolates there; every
  output pixel therefore gets a value.

  (b) TRUE --- nearest neighbour only copies existing values, so no labels that do not exist can
  appear. Bilinear or bicubic interpolation would create spurious intermediate label values at
  the boundaries.

  (c) FALSE --- bilinear interpolation yields a piecewise-linear function: continuous, but not
  differentiable at the grid lines. Bicubic interpolation is the one that is both continuous and
  differentiable.

  (d) FALSE --- the 8 corner voxels are what \emph{trilinear} interpolation uses. Keys' kernel
  has support $[-2,2]$, covering 4 samples per dimension, i.e.\ $4\times4 = 16$ pixels in 2D and
  $4\times4\times4 = 64$ voxels in 3D.
\end{solution}

% =============================================================== Question 5 ===
\question
Mixture models and Expectation-Maximization (EM) segmentation are essential building blocks for
many segmentation algorithms. Consider a multi-modal image in which every pixel carries
\emph{two} intensity values (for instance a co-registered T1-weighted and T2-weighted pair). How
many parameters does a Gaussian Mixture Model (GMM) with 3 components have? Count, for each
component, the mean vector, the full covariance matrix and the mixture weight.

\begin{solution}
  Answer: $3 \times (2 + 3 + 1) = 18$.

  Per component: a mean vector in $\mathbb{R}^2$ gives 2 parameters; a full $2\times2$ covariance
  matrix is symmetric and therefore has 3 free parameters ($\sigma_1^2$, $\sigma_2^2$ and one
  covariance); the mixture weight gives 1. That is 6 parameters per component and 18 in total.
\end{solution}
\answerspace{3.0cm}

% =============================================================== Question 6 ===
\question
An Active Shape Model (ASM) is fitted to an image iteratively. Once the model has been
initialised and candidate points have been found by sampling image features along the surface
normals of the model points, each iteration updates the model in two parts. Please briefly
describe what these two parts do, and state what keeps the fitted shape anatomically plausible.
You do not need to write down the equations.

\begin{solution}
  Answer: \emph{Pose estimation} computes a global transformation --- e.g.\ a similarity
  transform, i.e.\ translation, rotation and isotropic scaling --- that best aligns the model
  points with the candidate points found along the normals, for instance by minimising the mean
  Euclidean distance between the two point sets (ICP). \emph{Shape estimation} then removes the
  pose and computes the shape parameters $b$ of the PCA (point distribution) model in closed
  form, which deforms the mean shape along its modes of variation. The two are iterated until
  convergence.

  Plausibility is enforced by the model itself: only the leading principal modes are kept (e.g.\
  as many as needed to explain about 90\% of the variance), and each shape parameter is
  constrained to $|b_i| < 3\sigma_i$, so the fit can only produce shapes that are typical of the
  training population.
\end{solution}
\answerspace{7.5cm}

% =============================================================== Question 7 ===
\question
Thin-plate splines (TPS) and free-form deformations (FFD) with cubic B-splines are both
interpolation-based models for non-linear transformations. Which of the following statements on
their similarities and differences is \textbf{inaccurate}?

\begin{choices}
  \choice Both models place displacement vectors on a set of control points and interpolate in
  between with a smooth basis function, so that the resulting transformation is smooth for any
  value of the parameters.
  \choice The influence of a single control point is global in TPS, whereas for FFD with cubic
  B-splines it is local, because the B-spline basis functions have finite support.
  \choice In FFD, increasing the number of control points, i.e.\ using a finer control grid,
  results in smoother transformations.
  \choice For a given set of $N$ control points in 3D, both models are parameterised by the
  $3N$ coefficients of the displacement vectors at the control points.
\end{choices}

\begin{solution}
  Answer: C.

  It is the other way around: a finer control grid gives the model more degrees of freedom
  locally and therefore produces \emph{less} smooth transformations. This is precisely what makes
  multi-resolution registration possible --- start with few control points for a coarse, very
  smooth alignment and progressively refine the grid.
\end{solution}

% =============================================================== Question 8 ===
\question
Multi-atlas and patch-based methods were the state of the art in segmentation for a long time.
Which of the following statements about them is \textbf{incorrect}?

\begin{choices}
  \choice In multi-atlas segmentation, every atlas is registered to the unseen image, each atlas
  produces its own propagated segmentation, and these are combined in a decision fusion step such
  as majority voting or STAPLE.
  \choice STAPLE estimates the most probable underlying true segmentation and, simultaneously,
  the sensitivity and specificity of each of the input segmentations; it is solved iteratively
  with an Expectation-Maximization algorithm.
  \choice Selecting only the atlases that are similar to the unseen image before fusing performs
  worse than naively fusing all available atlases, because discarding atlases throws away
  information.
  \choice Patch-based label fusion requires only a coarse, rigid or affine registration, because
  correspondence is re-established locally by comparing patches inside a search window.
\end{choices}

\begin{solution}
  Answer: C.

  Atlas selection followed by fusion \emph{outperforms} naive fusion of all atlases: atlases that
  are dissimilar to the target contribute mostly registration errors, and removing them improves
  the fused result. The same reasoning underlies weighted voting, where atlases are weighted by
  how well they match the target rather than being discarded outright.
\end{solution}

% =============================================================== Question 9 ===
\question
Convolutional neural networks and Vision Transformers (ViT) build very different assumptions
into the model. Which of the following statements is \textbf{incorrect}?

\begin{choices}
  \choice In CNNs, locality and the two-dimensional neighbourhood structure of the image act as
  strong inductive biases for vision tasks.
  \choice In a ViT, self-attention is global and only the MLP layers act locally, so that
  splitting the image into patches is essentially the only two-dimensional inductive bias built
  into the model.
  \choice Because a ViT builds in fewer assumptions about images, it typically requires
  \emph{less} training data than a comparable CNN to reach the same accuracy.
  \choice Without positional encodings, a transformer encoder is equivariant to permutations of
  its input tokens, which is why positional information has to be added to the token embeddings.
\end{choices}

\begin{solution}
  Answer: C.

  Fewer built-in assumptions means the geometric properties of images that a CNN gets for free
  have to be learned from the data instead. ViTs therefore generally require \emph{more} training
  data than a comparable CNN, not less.
\end{solution}

% ============================================================== Question 10 ===
\question
You are training a U-Net to segment a small lesion that occupies far less than 1\% of the image.
Which of the following statements is correct?

\begin{choices}
  \choice With pixel-wise cross-entropy, the loss contributed by the lesion is typically much
  smaller than the loss contributed by the rest of the image, so the lesion can be dominated
  during training.
  \choice The S\o{}rensen-Dice coefficient computed on the binary label sets is differentiable
  and can therefore be used directly as a training loss.
  \choice The Dice loss makes the contribution of a class proportional to its size, which is why
  it helps for small structures.
  \choice The skip connections of the U-Net remove any need to reconsider the loss function,
  because they restore the resolution lost in the encoder.
  \choice Multi-class cross-entropy cannot be used when more than two classes are present.
\end{choices}

\begin{solution}
  Answer: A.

  The cross-entropy loss decomposes into a sum over pixels,
  $\mathcal{L} = \mathcal{L}_{\Omega_{small}} + \mathcal{L}_{\overline{\Omega_{small}}}$, and
  since $|\Omega_{small}| \ll |\overline{\Omega_{small}}|$ the first term is normally negligible.
  (B) is wrong because the Dice coefficient as a set measure is \emph{not} differentiable, which
  is why the soft Dice loss is used instead. (C) is wrong in the opposite direction: the soft Dice
  loss gives each class a contribution of comparable magnitude \emph{regardless} of its size,
  which is exactly why it helps here. (D) confuses resolution with class imbalance --- skip
  connections address the former, not the latter.
\end{solution}

% ============================================================== Question 11 ===
\question
Consider the encoder of a segmentation network that takes grey-scale input images of size
$128 \times 128$ with one input channel and has the following architecture:

\begin{itemize}
  \item \textbf{Conv1}: 32 kernels with kernel size $K = 3 \times 3$, stride $S = 1 \times 1$ in
    each direction, and ``same'' padding $P = 1$.
  \item \textbf{Pool1}: Max pooling with $K = 2 \times 2$ and $S = 2$.
  \item \textbf{Conv2}: 64 kernels with $K = 3 \times 3$, $S = 1$ and ``same'' padding $P = 1$.
  \item \textbf{Pool2}: Max pooling with $K = 2 \times 2$ and $S = 2$.
  \item \textbf{Conv3}: 128 kernels with $K = 3 \times 3$, $S = 1$ and ``same'' padding $P = 1$.
  \item \textbf{Pool3}: Max pooling with $K = 2 \times 2$ and $S = 2$.
\end{itemize}

\noindent Circle True or False:

\begin{parts}
  \part \TFT\ \ The global receptive field of a neuron in Pool3 is $22 \times 22$.

  \part \TFF\ \ The global stride of a neuron in Pool3 is $3 \times 3$, one for each pooling
  layer.

  \part \TFT\ \ The total number of convolutional kernel weights (excluding the biases) in
  Conv1, Pool1, Conv2, Pool2, Conv3 and Pool3 is
  $3 \cdot 3 \cdot 1 \cdot 32 + 3 \cdot 3 \cdot 32 \cdot 64 + 3 \cdot 3 \cdot 64 \cdot 128
  = 92\,448$.

  \part \TFF\ \ The output of Pool3 has a spatial size of $16 \times 16$, and each of its
  neurons therefore sees the entire $128 \times 128$ input image.
\end{parts}

\begin{solution}
  (a) TRUE --- track the receptive field $R$ and the jump (global stride) $j$ layer by layer,
  starting from $R = 1$, $j = 1$, using $R \leftarrow R + (K-1)\cdot j$ and $j \leftarrow j \cdot S$:
  Conv1 $\rightarrow R = 3,\ j = 1$; Pool1 $\rightarrow R = 4,\ j = 2$;
  Conv2 $\rightarrow R = 8,\ j = 2$; Pool2 $\rightarrow R = 10,\ j = 4$;
  Conv3 $\rightarrow R = 18,\ j = 4$; Pool3 $\rightarrow R = 22,\ j = 8$.

  (b) FALSE --- strides compound multiplicatively, they do not add: three poolings of stride 2
  give a global stride of $2 \cdot 2 \cdot 2 = 8$, i.e.\ $8 \times 8$.

  (c) TRUE --- $288 + 18\,432 + 73\,728 = 92\,448$. Pooling layers have no learnable parameters
  at all.

  (d) FALSE --- the spatial size is indeed $128 \to 64 \to 32 \to 16$, since ``same'' convolutions
  do not change the size and each pooling halves it. But the output size and the receptive field
  are different things: by (a) a neuron in Pool3 only sees a $22 \times 22$ window of the input,
  far less than the whole image. This is precisely why deeper networks, or skip connections
  combined with more downsampling, are needed to gather global context.
\end{solution}

% ============================================================== Question 12 ===
\question
Which of the following statements on uncertainty estimation, domain shift and learning from
unlabelled examples are correct (multiple answers are possible)?

\begin{choices}
  \choice Standard, deterministic training and inference amount to approximating
  $p(\theta \mid D, M)$ by a Dirac delta at $\theta^*$ and $p(y \mid x, \theta, M)$ by a Dirac
  delta at $f(x;\theta)$, so that the resulting predictive distribution is a point mass carrying
  no uncertainty at all.
  \choice Test-time adaptation still requires labelled samples from the target domain, but fewer
  of them than fine-tuning does.
  \choice Self-training with pseudo-labels works well only if the predictions of the initial
  model are reasonably good; with poor initial estimates it can diverge quickly.
  \choice In the mean-teacher model the teacher is trained by minimising the supervised loss on
  the labelled examples, while the student is the exponential moving average of the teacher.
  \choice Domain generalisation by extreme data augmentation trains only in the source domain,
  but the improvement it brings may be limited to the kinds of variation that the augmentations
  actually cover.
\end{choices}

\begin{solution}
  Answer: (A), (C) and (E) are correct.

  (B) is wrong: test-time adaptation uses \emph{no} labels and no training set at the target
  domain --- it adapts a subset of the parameters at inference time by minimising a label-free
  cost (prediction entropy, an auto-encoder reconstruction loss, or the agreement of the output
  with a source-trained segmentation prior). It is fine-tuning that needs a few labelled target
  samples.

  (D) has the two networks swapped: the student is the one trained on the supervised plus
  consistency loss, and the teacher is the exponential moving average of the student,
  $\phi^i = \alpha \phi^{i-1} + (1-\alpha)\theta^i$.
\end{solution}

\end{questions}

\end{document}
"
%% Compile the blank version (no answers):  pdflatex Sample_Exam_2.tex
%%
\ifx\withanswers\undefined
  \documentclass[11pt]{exam}
\else
  \documentclass[11pt,answers]{exam}
\fi

\usepackage[margin=2.4cm]{geometry}
\usepackage{amsmath,amssymb}
\usepackage{array}
\usepackage[T1]{fontenc}

% ---------------------------------------------------------------- page style
\pagestyle{foot}
\firstpageheader{}{}{}
\firstpagefooter{}{}{}
\runningheader{}{}{}
\runningfooter{}{}{Page \thepage}

% -------------------------------------------------------------- multiplechoice
% Reproduce the "(A) (B) (C) ..." labelling of the original exam.
\renewcommand{\choicelabel}{$\bigcirc$~(\Alph{choice})}

% ---------------------------------------------------------- true/false macros
% \TFT -> the statement is TRUE, \TFF -> the statement is FALSE.
% In the answer version the correct word is set in capitals and bold.
\newcommand{\TFT}{\ifprintanswers\textbf{TRUE}\hspace{1.1em}False\else True\hspace{1.1em}False\fi}
\newcommand{\TFF}{True\hspace{1.1em}\ifprintanswers\textbf{FALSE}\else False\fi}

% ------------------------------------------------------------- writing space
% Reserves room to write in the blank version; collapses in the answer version.
\newcommand{\answerspace}[1]{\ifprintanswers\else\vspace{#1}\fi}

\begin{document}

% ------------------------------------------------------------------ title page
\begin{center}
  {\LARGE\bfseries Sample Exam}\\[2mm]
  {\Large Medical Image Analysis}\\[6mm]
\end{center}

\vspace{4mm}
\noindent Student name:\hrulefill

\vspace{6mm}
\noindent Immatriculation number:\hrulefill

\vspace{10mm}

\begin{center}
  {\large\bfseries Exam Questions}
\end{center}

\noindent\textbf{Note:} All the questions have equal weight.

\vspace{4mm}

\begin{questions}

% =============================================================== Question 1 ===
\question
In 2D acquisitions, 2D image slices are acquired in one pre-defined direction one after the
other, and the final 3D volume is created by stacking the 2D slices. For instance, in axial
acquisitions 2D slices are acquired in the transverse plane, in sagittal acquisitions in the
sagittal plane, and in coronal acquisitions in the coronal plane. The voxel sizes of 2D
acquisitions are given based on the acquisition direction, as in-plane and through-plane voxel
sizes, where the first is represented with two values (pixel size in the imaging plane) and the
latter corresponds to a single value (length of the voxel in the third dimension). Recall that
the field-of-view is the real-world size of the entire image, i.e.\ image size $\times$ voxel
size. Based on this please compute the following.

\begin{parts}

  \part A \emph{sagittal} acquisition with voxel size in-plane $0.5 \times 0.5$~mm$^2$ with 384
  pixels in both directions and voxel size through-plane $4.0$~mm with 60 pixels is acquired.
  What is the field-of-view in the head-to-toe, left-to-right and anterior-to-posterior
  directions, respectively?

  \begin{solution}
    Answer: $192 \times 240 \times 192$~mm$^3$.

    The sagittal plane contains the head-to-toe and the anterior-to-posterior directions, so
    both in-plane directions give $384 \times 0.5 = 192$~mm. The through-plane direction of a
    sagittal acquisition is left-to-right: $60 \times 4.0 = 240$~mm.
  \end{solution}
\answerspace{2.2cm}

  \part One year later the same patient is scanned again, this time with an \emph{axial}
  acquisition with voxel size in-plane $0.8 \times 0.8$~mm$^2$ with a field-of-view of 192~mm in
  both directions, and voxel size through-plane of $2.0$~mm with a field-of-view of 240~mm.
  What are the number of voxels in the head-to-toe, left-to-right and anterior-to-posterior
  directions, respectively?

  \begin{solution}
    Answer: $120 \times 240 \times 240$~voxels.

    The through-plane direction of an axial acquisition is head-to-toe: $240 / 2.0 = 120$. Both
    in-plane directions give $192 / 0.8 = 240$.
  \end{solution}
\answerspace{2.2cm}

  \part You want to compare a volumetric measurement of the same anatomical structure in the two
  images. Following the guideline given in the lecture, you resample the second image (b) onto
  the voxel grid of the first image (a). What are the resize factors along the head-to-toe,
  left-to-right and anterior-to-posterior directions, and what is the resulting image size?
  Why do you match the voxel size and not the image size?

  \begin{solution}
    Answer: resize factors $4.0 \times 0.2 \times 1.6$, resulting image size
    $480 \times 48 \times 384$~voxels.

    The resize factor is the ratio of the old to the new voxel size, per direction:
    head-to-toe $2.0/0.5 = 4.0$, left-to-right $0.8/4.0 = 0.2$, anterior-to-posterior
    $0.8/0.5 = 1.6$. Applying these to (b): $120 \cdot 4.0 = 480$, $240 \cdot 0.2 = 48$,
    $240 \cdot 1.6 = 384$. The field-of-view is unchanged ($240 \times 192 \times 192$~mm$^3$),
    as it must be.

    Measurements taken in an image correspond to real-world quantities only through the voxel
    size, so the voxel size is what has to be matched. Matching the image sizes instead would
    silently change the effective voxel size (and here also the aspect ratio), and the two
    measurements would no longer be comparable.
  \end{solution}
\answerspace{4.5cm}

\end{parts}

% =============================================================== Question 2 ===
\question
Many noise suppression methods can be written as a weighted average
$\hat I(x) = \sum_{y} w(x,y)\, J(y)$ of the noisy image $J$, with $w(x,y) \geq 0$ and
$\sum_y w(x,y) = 1$. Which of the two weight definitions below is the one used by the
\emph{non-local means} filter? $N(x)$ denotes the image patch centred at $x$, $Z(x)$ is a
normalisation constant and $h$ the degree of filtering. Please circle the correct choice.

\vspace{2mm}
\begin{center}
  $\displaystyle w(x,y) = \frac{1}{Z(x)} \exp\left(-\frac{\|J(x) - J(y)\|_2^2}{h^2}\right)$
  \hspace{1.6cm}
  $\displaystyle w(x,y) = \frac{1}{Z(x)} \exp\left(-\frac{\|J(N(x)) - J(N(y))\|_2^2}{h^2}\right)$
\end{center}
\vspace{2mm}

\begin{solution}
  Answer: Right one.

  Non-local means compares the whole geometric configuration of a \emph{patch}, not a single
  grey value. The left expression compares single (noisy) intensities and is the
  Yaroslavsky / bilateral type of neighbourhood filter, which is much less robust: two pixels
  can share an intensity by chance while sitting in completely different structures.
\end{solution}
\answerspace{1.2cm}

% =============================================================== Question 3 ===
\question
Deformable intensity-based registration is usually written in one of the two equivalent forms
below. The first is derived from a probabilistic model, the second is energy-based;
$T_\theta$ is the transformation with parameters $\theta$, $I$ the target and $J$ the moving
image.

\vspace{2mm}
\begin{center}
  $\displaystyle \max_\theta\ \log p(I \mid T_\theta \circ J) + \log p(T_\theta)$
  \hspace{1.8cm}
  $\displaystyle \min_\theta\ d(I, T_\theta \circ J) + \lambda\, r(T_\theta^{-1})$
\end{center}
\vspace{2mm}

\begin{parts}
  \part Please indicate which term corresponds to which term in the two formulations.

  \begin{solution}
    Answer: $\log p(I \mid T_\theta \circ J) \Longleftrightarrow d(I, T_\theta \circ J)$ and
    $\log p(T_\theta) \Longleftrightarrow \lambda\, r(T_\theta^{-1})$.
  \end{solution}
\answerspace{2.2cm}

  \part Deformable registration is an ill-posed problem. Which of the two terms addresses this,
  and how?

  \begin{solution}
    Answer: the regularisation term $\lambda\, r(T_\theta^{-1})$, equivalently the prior
    $p(T_\theta)$.

    Many very different transformations explain the intensities about equally well, so the data
    term alone does not determine a unique solution. The regulariser expresses a preference among
    them --- e.g.\ $\ell_2$ regularisation of $\nabla u$ prefers smoothly varying displacement
    fields, $\ell_1$ prefers piecewise constant ones, and elastic / hyperelastic energies prefer
    small or locally rigid deformations.
  \end{solution}
\answerspace{3.5cm}
\end{parts}

% =============================================================== Question 4 ===
\question
Please indicate whether the statement is true or false by circling the correct choice in each of
the following.

\begin{parts}
  \part \TFF\ \ In the discrete setting, the \emph{forward} transformation is preferred over
  the backward (inverse) transformation for resampling an image, because it guarantees that
  every pixel of the output grid is assigned a value.

  \part \TFT\ \ Nearest-neighbour interpolation is the appropriate choice for transforming a
  label (segmentation) map, because it does not introduce new values.

  \part \TFF\ \ Bilinear interpolation produces a function that is both continuous and
  differentiable.

  \part \TFF\ \ In 3D, Keys' cubic convolution (tricubic) interpolation uses only the 8 corner
  voxels of the cell that contains the query point.
\end{parts}

\begin{solution}
  (a) FALSE --- it is exactly the other way around. Under a forward mapping the transformed
  input pixels do not land on output grid points, two input pixels may land on the same output
  pixel, and some output pixels receive nothing at all (holes). The backward mapping asks, for
  each output pixel, where it comes from in the input image, and interpolates there; every
  output pixel therefore gets a value.

  (b) TRUE --- nearest neighbour only copies existing values, so no labels that do not exist can
  appear. Bilinear or bicubic interpolation would create spurious intermediate label values at
  the boundaries.

  (c) FALSE --- bilinear interpolation yields a piecewise-linear function: continuous, but not
  differentiable at the grid lines. Bicubic interpolation is the one that is both continuous and
  differentiable.

  (d) FALSE --- the 8 corner voxels are what \emph{trilinear} interpolation uses. Keys' kernel
  has support $[-2,2]$, covering 4 samples per dimension, i.e.\ $4\times4 = 16$ pixels in 2D and
  $4\times4\times4 = 64$ voxels in 3D.
\end{solution}

% =============================================================== Question 5 ===
\question
Mixture models and Expectation-Maximization (EM) segmentation are essential building blocks for
many segmentation algorithms. Consider a multi-modal image in which every pixel carries
\emph{two} intensity values (for instance a co-registered T1-weighted and T2-weighted pair). How
many parameters does a Gaussian Mixture Model (GMM) with 3 components have? Count, for each
component, the mean vector, the full covariance matrix and the mixture weight.

\begin{solution}
  Answer: $3 \times (2 + 3 + 1) = 18$.

  Per component: a mean vector in $\mathbb{R}^2$ gives 2 parameters; a full $2\times2$ covariance
  matrix is symmetric and therefore has 3 free parameters ($\sigma_1^2$, $\sigma_2^2$ and one
  covariance); the mixture weight gives 1. That is 6 parameters per component and 18 in total.
\end{solution}
\answerspace{3.0cm}

% =============================================================== Question 6 ===
\question
An Active Shape Model (ASM) is fitted to an image iteratively. Once the model has been
initialised and candidate points have been found by sampling image features along the surface
normals of the model points, each iteration updates the model in two parts. Please briefly
describe what these two parts do, and state what keeps the fitted shape anatomically plausible.
You do not need to write down the equations.

\begin{solution}
  Answer: \emph{Pose estimation} computes a global transformation --- e.g.\ a similarity
  transform, i.e.\ translation, rotation and isotropic scaling --- that best aligns the model
  points with the candidate points found along the normals, for instance by minimising the mean
  Euclidean distance between the two point sets (ICP). \emph{Shape estimation} then removes the
  pose and computes the shape parameters $b$ of the PCA (point distribution) model in closed
  form, which deforms the mean shape along its modes of variation. The two are iterated until
  convergence.

  Plausibility is enforced by the model itself: only the leading principal modes are kept (e.g.\
  as many as needed to explain about 90\% of the variance), and each shape parameter is
  constrained to $|b_i| < 3\sigma_i$, so the fit can only produce shapes that are typical of the
  training population.
\end{solution}
\answerspace{7.5cm}

% =============================================================== Question 7 ===
\question
Thin-plate splines (TPS) and free-form deformations (FFD) with cubic B-splines are both
interpolation-based models for non-linear transformations. Which of the following statements on
their similarities and differences is \textbf{inaccurate}?

\begin{choices}
  \choice Both models place displacement vectors on a set of control points and interpolate in
  between with a smooth basis function, so that the resulting transformation is smooth for any
  value of the parameters.
  \choice The influence of a single control point is global in TPS, whereas for FFD with cubic
  B-splines it is local, because the B-spline basis functions have finite support.
  \choice In FFD, increasing the number of control points, i.e.\ using a finer control grid,
  results in smoother transformations.
  \choice For a given set of $N$ control points in 3D, both models are parameterised by the
  $3N$ coefficients of the displacement vectors at the control points.
\end{choices}

\begin{solution}
  Answer: C.

  It is the other way around: a finer control grid gives the model more degrees of freedom
  locally and therefore produces \emph{less} smooth transformations. This is precisely what makes
  multi-resolution registration possible --- start with few control points for a coarse, very
  smooth alignment and progressively refine the grid.
\end{solution}

% =============================================================== Question 8 ===
\question
Multi-atlas and patch-based methods were the state of the art in segmentation for a long time.
Which of the following statements about them is \textbf{incorrect}?

\begin{choices}
  \choice In multi-atlas segmentation, every atlas is registered to the unseen image, each atlas
  produces its own propagated segmentation, and these are combined in a decision fusion step such
  as majority voting or STAPLE.
  \choice STAPLE estimates the most probable underlying true segmentation and, simultaneously,
  the sensitivity and specificity of each of the input segmentations; it is solved iteratively
  with an Expectation-Maximization algorithm.
  \choice Selecting only the atlases that are similar to the unseen image before fusing performs
  worse than naively fusing all available atlases, because discarding atlases throws away
  information.
  \choice Patch-based label fusion requires only a coarse, rigid or affine registration, because
  correspondence is re-established locally by comparing patches inside a search window.
\end{choices}

\begin{solution}
  Answer: C.

  Atlas selection followed by fusion \emph{outperforms} naive fusion of all atlases: atlases that
  are dissimilar to the target contribute mostly registration errors, and removing them improves
  the fused result. The same reasoning underlies weighted voting, where atlases are weighted by
  how well they match the target rather than being discarded outright.
\end{solution}

% =============================================================== Question 9 ===
\question
Convolutional neural networks and Vision Transformers (ViT) build very different assumptions
into the model. Which of the following statements is \textbf{incorrect}?

\begin{choices}
  \choice In CNNs, locality and the two-dimensional neighbourhood structure of the image act as
  strong inductive biases for vision tasks.
  \choice In a ViT, self-attention is global and only the MLP layers act locally, so that
  splitting the image into patches is essentially the only two-dimensional inductive bias built
  into the model.
  \choice Because a ViT builds in fewer assumptions about images, it typically requires
  \emph{less} training data than a comparable CNN to reach the same accuracy.
  \choice Without positional encodings, a transformer encoder is equivariant to permutations of
  its input tokens, which is why positional information has to be added to the token embeddings.
\end{choices}

\begin{solution}
  Answer: C.

  Fewer built-in assumptions means the geometric properties of images that a CNN gets for free
  have to be learned from the data instead. ViTs therefore generally require \emph{more} training
  data than a comparable CNN, not less.
\end{solution}

% ============================================================== Question 10 ===
\question
You are training a U-Net to segment a small lesion that occupies far less than 1\% of the image.
Which of the following statements is correct?

\begin{choices}
  \choice With pixel-wise cross-entropy, the loss contributed by the lesion is typically much
  smaller than the loss contributed by the rest of the image, so the lesion can be dominated
  during training.
  \choice The S\o{}rensen-Dice coefficient computed on the binary label sets is differentiable
  and can therefore be used directly as a training loss.
  \choice The Dice loss makes the contribution of a class proportional to its size, which is why
  it helps for small structures.
  \choice The skip connections of the U-Net remove any need to reconsider the loss function,
  because they restore the resolution lost in the encoder.
  \choice Multi-class cross-entropy cannot be used when more than two classes are present.
\end{choices}

\begin{solution}
  Answer: A.

  The cross-entropy loss decomposes into a sum over pixels,
  $\mathcal{L} = \mathcal{L}_{\Omega_{small}} + \mathcal{L}_{\overline{\Omega_{small}}}$, and
  since $|\Omega_{small}| \ll |\overline{\Omega_{small}}|$ the first term is normally negligible.
  (B) is wrong because the Dice coefficient as a set measure is \emph{not} differentiable, which
  is why the soft Dice loss is used instead. (C) is wrong in the opposite direction: the soft Dice
  loss gives each class a contribution of comparable magnitude \emph{regardless} of its size,
  which is exactly why it helps here. (D) confuses resolution with class imbalance --- skip
  connections address the former, not the latter.
\end{solution}

% ============================================================== Question 11 ===
\question
Consider the encoder of a segmentation network that takes grey-scale input images of size
$128 \times 128$ with one input channel and has the following architecture:

\begin{itemize}
  \item \textbf{Conv1}: 32 kernels with kernel size $K = 3 \times 3$, stride $S = 1 \times 1$ in
    each direction, and ``same'' padding $P = 1$.
  \item \textbf{Pool1}: Max pooling with $K = 2 \times 2$ and $S = 2$.
  \item \textbf{Conv2}: 64 kernels with $K = 3 \times 3$, $S = 1$ and ``same'' padding $P = 1$.
  \item \textbf{Pool2}: Max pooling with $K = 2 \times 2$ and $S = 2$.
  \item \textbf{Conv3}: 128 kernels with $K = 3 \times 3$, $S = 1$ and ``same'' padding $P = 1$.
  \item \textbf{Pool3}: Max pooling with $K = 2 \times 2$ and $S = 2$.
\end{itemize}

\noindent Circle True or False:

\begin{parts}
  \part \TFT\ \ The global receptive field of a neuron in Pool3 is $22 \times 22$.

  \part \TFF\ \ The global stride of a neuron in Pool3 is $3 \times 3$, one for each pooling
  layer.

  \part \TFT\ \ The total number of convolutional kernel weights (excluding the biases) in
  Conv1, Pool1, Conv2, Pool2, Conv3 and Pool3 is
  $3 \cdot 3 \cdot 1 \cdot 32 + 3 \cdot 3 \cdot 32 \cdot 64 + 3 \cdot 3 \cdot 64 \cdot 128
  = 92\,448$.

  \part \TFF\ \ The output of Pool3 has a spatial size of $16 \times 16$, and each of its
  neurons therefore sees the entire $128 \times 128$ input image.
\end{parts}

\begin{solution}
  (a) TRUE --- track the receptive field $R$ and the jump (global stride) $j$ layer by layer,
  starting from $R = 1$, $j = 1$, using $R \leftarrow R + (K-1)\cdot j$ and $j \leftarrow j \cdot S$:
  Conv1 $\rightarrow R = 3,\ j = 1$; Pool1 $\rightarrow R = 4,\ j = 2$;
  Conv2 $\rightarrow R = 8,\ j = 2$; Pool2 $\rightarrow R = 10,\ j = 4$;
  Conv3 $\rightarrow R = 18,\ j = 4$; Pool3 $\rightarrow R = 22,\ j = 8$.

  (b) FALSE --- strides compound multiplicatively, they do not add: three poolings of stride 2
  give a global stride of $2 \cdot 2 \cdot 2 = 8$, i.e.\ $8 \times 8$.

  (c) TRUE --- $288 + 18\,432 + 73\,728 = 92\,448$. Pooling layers have no learnable parameters
  at all.

  (d) FALSE --- the spatial size is indeed $128 \to 64 \to 32 \to 16$, since ``same'' convolutions
  do not change the size and each pooling halves it. But the output size and the receptive field
  are different things: by (a) a neuron in Pool3 only sees a $22 \times 22$ window of the input,
  far less than the whole image. This is precisely why deeper networks, or skip connections
  combined with more downsampling, are needed to gather global context.
\end{solution}

% ============================================================== Question 12 ===
\question
Which of the following statements on uncertainty estimation, domain shift and learning from
unlabelled examples are correct (multiple answers are possible)?

\begin{choices}
  \choice Standard, deterministic training and inference amount to approximating
  $p(\theta \mid D, M)$ by a Dirac delta at $\theta^*$ and $p(y \mid x, \theta, M)$ by a Dirac
  delta at $f(x;\theta)$, so that the resulting predictive distribution is a point mass carrying
  no uncertainty at all.
  \choice Test-time adaptation still requires labelled samples from the target domain, but fewer
  of them than fine-tuning does.
  \choice Self-training with pseudo-labels works well only if the predictions of the initial
  model are reasonably good; with poor initial estimates it can diverge quickly.
  \choice In the mean-teacher model the teacher is trained by minimising the supervised loss on
  the labelled examples, while the student is the exponential moving average of the teacher.
  \choice Domain generalisation by extreme data augmentation trains only in the source domain,
  but the improvement it brings may be limited to the kinds of variation that the augmentations
  actually cover.
\end{choices}

\begin{solution}
  Answer: (A), (C) and (E) are correct.

  (B) is wrong: test-time adaptation uses \emph{no} labels and no training set at the target
  domain --- it adapts a subset of the parameters at inference time by minimising a label-free
  cost (prediction entropy, an auto-encoder reconstruction loss, or the agreement of the output
  with a source-trained segmentation prior). It is fine-tuning that needs a few labelled target
  samples.

  (D) has the two networks swapped: the student is the one trained on the supervised plus
  consistency loss, and the teacher is the exponential moving average of the student,
  $\phi^i = \alpha \phi^{i-1} + (1-\alpha)\theta^i$.
\end{solution}

\end{questions}

\end{document}
